\section{The sequential input}\label{app:sequential}
We recall what is used from the companion paper \cite{Companion} on the sequential chain
\(S_0,\dots,S_n\). Its Theorem~1.1 states that every compatibility class contains at most
\(n\) sinks (nonhyperbolic ones included), and that for every \(n\ge1\) there are positive
rational rate constants and totals with exactly \(2n-1\) positive equilibria, all hyperbolic,
of which \(n\) are sinks and \(n-1\) are saddles with a one-dimensional unstable manifold. The
bound \(2n-1\) on the number of equilibria is due to Wang and Sontag \cite{WangSontag2008}.

The construction is as follows. For auxiliary numbers \(x_0,\dots,x_{2n-2}>1\) a positive
polynomial recurrence produces a chain whose \(2n-1\) equilibria, at the totals
\((2r,2,2r+2)\), have the free enzyme concentrations
\begin{equation}\label{eq:EFx}
 E=\frac{x_j-1}{2},\qquad F=\frac{4}{x_j+1},\qquad u=\frac{x_j^2-1}{8}.
\end{equation}
Positivity of all coefficients holds whenever \(r>\max_ju(x_j)\) and
\(4r+\sqrt{1+8r}\ge\sum_jx_j\), by an interlacing argument. The loadings and effective rates
are determined by the equilibrium data, while \(n\) catalytic scales and the dissociation
constants remain free. Stability is obtained at the \emph{coalesced} system \(x_j\equiv3\),
where all equilibria coincide at a point \(X_\circ\) with \(E=F=1\): after an acceleration
of binding as in \eqref{eq:loading}, the slow dynamics on the retained pools is governed by an
\(n\times n\) matrix whose leading \((n-1)\)-block converges, as \(r\to\infty\), to a path
matrix with a rank-one positive feedback of gain strictly less than one, hence is Hurwitz, and
a small final catalytic current keeps the remaining eigenvalue at zero. The outcome is
\cite[Lemma~7.2]{Companion}: for suitable positive constants the class Jacobian at \(X_\circ\)
has an algebraically simple eigenvalue \(0\) and \(3n-1\) eigenvalues in the open left
half-plane. With these constants frozen, the \emph{split family} takes
\(x_j=3+\delta\xi_j\) for real \(\xi_0<\dots<\xi_{2n-2}\). All data are real-analytic in
\(\delta\), the totals are fixed, and \cite[Theorem~1.2]{Companion} gives the splitting law
\eqref{eq:split}, with the constant \(a>0\) of \cite[(7.2)]{Companion}, from an exact formula for
the determinant of the class Jacobian. Since \(\prod_{k\ne j}(\xi_j-\xi_k)\) has the sign
\((-1)^j\), the even-indexed equilibria are the sinks. By \eqref{eq:EFx}, \(E=1+\delta\xi_j/2\)
at \(X_j(\delta)\). For \(n=1\) the equilibrium is unique and globally attracting
\cite{AngeliSontag2008}, and \cite{Companion} exhibits rate constants for which it is
hyperbolic.

\section{Proof of Theorem~\ref{thm:action}}\label{app:action}

\subsection{Phase and fibre coordinates}
Extend the split family by the equation \(\dot\delta=0\). At \((X_\circ,0)\) the extended
linearization has a two-dimensional centre space, on which it vanishes because the zero
eigenvalue is simple and \(\delta\) is a parameter, and a Hurwitz complement. For every \(k\)
there is a \(C^k\) centre manifold; it is normally attracting and contains all equilibria
\(X_j(\delta)\) for small \(\delta\) \cite{Carr1981}. On a small neighbourhood the flow on the
centre manifold and its inverse grow at most like \(e^{\varepsilon_0|t|}\), while the stable
directions contract at a rate \(\beta_0>0\), and \((k+1)\varepsilon_0<\beta_0\) for a small
enough neighbourhood. Under this spectral gap condition the stable foliation of the centre
manifold is \(C^k\) \cite{Fenichel1979,HirschPughShub1977,ElderingKvalheimRevzen2018}. If a
compact invariant base is required by the cited theorem, one modifies the nonlinear terms
outside the neighbourhood, which can be done with arbitrarily small \(C^1\) norm because they
vanish to first order, and restricts afterwards. Let \(\pi\) be the projection along the stable
fibres onto the centre manifold of parameter \(\delta\). We take
\[
 z=2\,(E\circ\pi)-2 .
\]
This is an admissible coordinate on the one-dimensional centre manifold because the centre
direction at \(X_\circ\) is tangent to the equilibrium curve \eqref{eq:EFx}, along which
\(\dd E/\dd x=1/2\ne0\); and \(z=\delta\xi_j\) at \(X_j(\delta)\) exactly. Choosing fibre
coordinates \(w\) that vanish on the centre manifold gives \eqref{eq:normalform} with
\(\dot z=f(z,\delta)\) \emph{independent of} \(w\), by invariance of the foliation, and with
\(\mathsf B(0,0,0)\) Hurwitz.

Since \(f(\delta\xi_j,\delta)=0\) for all \(j\), repeated division (Hadamard's lemma, with
\(k\) large) gives \(f(z,\delta)=g(z,\delta)\prod_j(z-\delta\xi_j)\) with \(g\) continuous.
The eigenvalue of the class Jacobian along the centre manifold at \(X_j(\delta)\) is
\(\partial_zf(\delta\xi_j,\delta)=g(\delta\xi_j,\delta)\,\delta^{2n-2}\prod_{k\ne j}(\xi_j-\xi_k)\),
and comparison with \eqref{eq:split} gives \(g(0,0)=-a\). This is a factorization of the
\emph{dynamics} on the centre manifold, not only a statement about its zero set.

\subsection{Lower bound for the action}
Put \(\omega(x)=D(z\circ\pi)(x)\), \(b(x)=\omega(x)Q(x)\omega(x)^T\) and
\(b_0=b(X_\circ)\). Near the interior point \(X_\circ\) all propensities are positive, and the
reaction vectors span the class, so \(Q\) is positive definite there and \(b_0>0\); moreover
the cone generated by the reaction vectors is the whole class, since every binding step has
its dissociation and a substrate transfer can be reversed through the opposite enzyme. Hence
\(\mathcal H(x,\cdot)\) in \eqref{eq:hamiltonian} is strictly convex and superlinear, and
\[
 \mathcal H(x,\theta)=H(x)\cdot\theta+\tfrac12\theta^TQ(x)\theta+O(|\theta|^3),\qquad
 \cI(x,H(x)+\chi)=\tfrac12\chi^TQ(x)^{-1}\chi+O(|\chi|^3)
\]
uniformly near \(X_\circ\). Let \(U(z)=-\int_{\delta\xi_i}^zf(y,\delta)\,\dd y\). Because
\(i\) is even, \(U\ge0\) between the cuts, and by \eqref{eq:normalform}
\begin{equation}\label{eq:Ucut}
 U(\delta s)=a\,\delta^{2n}\int_{\xi_i}^{s}\prod_j(y-\xi_j)\,\dd y\,(1+O(\delta)),
 \qquad s\in\{l,r_{\mathrm c}\}.
\end{equation}
Let \(W_1=2U(z\circ\pi)/b_\delta\) with \(b_\delta=b_0+C_1\delta\). Then
\(DW_1=-(2f/b_\delta)\,\omega\), and since \(\omega\cdot H=f\) \emph{exactly},
\[
 \mathcal H(x,DW_1)=-\frac{2f^2}{b_\delta}\Bigl(1-\frac{b(x)}{b_\delta}\Bigr)+O(|f|^3)\le0
\]
on \(D_\delta\) if \(C_1\) is large, because \(b(x)=b_0+O(\delta)\) and
\(|f|=O(\delta^{2n-1})\) there. The exact autonomy of the phase is essential for this
inequality at points with \(w\ne0\). For every path \(\phi\) the Fenchel inequality
\(\cI(\phi,\dot\phi)\ge DW_1(\phi)\cdot\dot\phi-\mathcal H(\phi,DW_1(\phi))\) gives an action
of at least the increase of \(W_1\); a path leaving \(D_\delta\) through a phase cut therefore
costs at least \(2U(\delta s)/b_\delta\). For exits through the transverse boundary use
\(W_2=\kappa_tw^TP_tw\). Its derivative along the flow is at most \(-c\kappa_t\|w\|^2\) for
small \(\delta\), and the quadratic and higher terms of \(\mathcal H(x,DW_2)\) are at most
\(C\kappa_t^2\|w\|^2\); with \(\kappa_t\) small, \(\mathcal H(x,DW_2)\le0\), and a transverse
exit costs at least \(\kappa_t\delta^{2\alpha_t}\), which exceeds \(\delta^{2n}\) in order
because \(\alpha_t<n\). Every exit from \(D_\delta\) is of one of the two kinds.

\subsection{A matching path}
Consider the controlled field \(H+\chi\) with
\(\chi(x)=-\bigl(2f(z,\delta)/b(x)\bigr)\,Q(x)\,\omega(x)^T\). Its phase satisfies
\(\dot z=f-2f=-f\) exactly, so the phase moves from near \(\delta\xi_i\) to either cut. The
fibre equation is forced by a term of order \(|f|=O(\delta^{2n-1})\), so
\(\|w\|=O(\delta^{2n-1})\) stays inside the transverse tube because \(2n-1>\alpha_t\). The cost
rate is \(\tfrac12\chi^TQ^{-1}\chi+O(|\chi|^3)=2f^2/b(x)+O(|f|^3)\), and with
\(\dd t=\dd z/(-f)\) the total cost is \(2U(\delta s)(1+O(\delta))/b_0\). The path is started at
a point arbitrarily close to the equilibrium, which it leaves in finite time; joining that
point to the equilibrium costs arbitrarily little because \(Q\) is positive definite. Together
with the lower bound and \eqref{eq:Ucut}, this proves \eqref{eq:action}.

\subsection{Exit times from the generator}
Fix \(\delta>0\), and let \(J_-=\min\{2U(\delta l)/b_\delta,\,2U(\delta r_{\mathrm c})/b_\delta,\,
\kappa_t\delta^{2\alpha_t}\}\) and \(J_+=I_\delta\).

\emph{Lower tail.} Enlarging \(C_1\) and decreasing \(\kappa_t\) slightly makes
\(\mathcal H(x,DW_j)\) strictly negative on \(\{W_j\ge\varepsilon\}\cap\overline{D_\delta}\),
for any fixed \(\varepsilon>0\). A Taylor expansion on the set of states reachable by one jump
from \(D_\delta\) gives
\(e^{-\Omega W_j}\cL_\Omega e^{\Omega W_j}=\Omega\,\mathcal H(x,DW_j)+O_\delta(1)\), which is
nonpositive on \(\{W_j\ge\varepsilon\}\) for large \(\Omega\) and at most
\(C_\delta\Omega\) elsewhere. Dynkin's formula for the stopped process and Markov's
inequality, applied to \(e^{\Omega W_1}\) and \(e^{\Omega W_2}\) separately, give for initial
states with \(W_j(X_0)\le\varepsilon\)
\begin{equation}\label{eq:exitlower}
 \Pr(\tau_\delta\le t)\le C'_\delta\,(1+C_\delta\Omega t)\,e^{-\Omega(J_--2\varepsilon)};
\end{equation}
the overshoot changes the boundary values of \(W_j\) by \(O_\delta(1/\Omega)\), which is
absorbed in \(C'_\delta\). Hence \(\Pr(\tau_\delta\le e^{\Omega(J_--4\varepsilon)})\to0\).

\emph{Upper tail.} The rate function has the flux representation
\[
 \cI(x,y)=\inf\Bigl\{\sum_\rho\bigl(q_\rho\log\frac{q_\rho}{a_\rho(x)}-q_\rho+a_\rho(x)\bigr):
 \ q_\rho\ge0,\ \sum_\rho q_\rho\nu_\rho=y\Bigr\}
\]
\cite{ShwartzWeiss1995}. Approximate the path of the previous subsection by a path on a finite
time interval with bounded positive fluxes \(q_\rho(t)\), cost at most \(J_++\varepsilon\), that
ends at a positive distance outside \(D_\delta\). Under the change of measure with intensities
\(\Omega q_\rho\), the compensated counts divided by \(\Omega\) have quadratic variation
\(O(1/\Omega)\), so by Doob's and Gronwall's inequalities the process follows the path with
probability tending to one, and the logarithm of the likelihood ratio, divided by \(\Omega\),
converges in probability to minus the cost, all logarithms being bounded on the compact
positive tube. Hence, from every starting point in a small neighbourhood \(G_0\) of the sink,
the original chain exits within the duration of the path with probability at least
\(\tfrac12e^{-\Omega(J_++2\varepsilon)}\) for large \(\Omega\). Every deterministic
trajectory from \(\overline{D_\delta}\) enters \(G_0\) within a common finite time, since the
phase field points inward and has the single zero \(\delta\xi_i\) between the cuts and the
fibres contract; by the same martingale estimate the chain either exits or reaches \(G_0\)
within that time with probability at least \(1/2\). The strong Markov property then gives, for
a block length \(h_\delta\),
\begin{equation}\label{eq:exitupper}
 \Pr(\tau_\delta>t)\le\bigl(1-\tfrac14e^{-\Omega(J_++2\varepsilon)}\bigr)^{\lfloor t/h_\delta\rfloor},
\end{equation}
which tends to zero for \(t=e^{\Omega(J_++3\varepsilon)}\). This direct argument avoids
applying a general exit theorem for unconstrained diffusions or jump processes to a process
confined to a conserved class; the general large-deviation theory of reaction networks has
hypotheses that would have to be verified \cite{AgazziDemboEckmann2018}.

\emph{Conclusion.} By \eqref{eq:action} and \eqref{eq:Ucut},
\(J_\pm=c_i\delta^{2n}(1+O(\delta))\). Given \(\zeta>0\), choose first \(\delta\) small, then
\(\varepsilon<\zeta\delta^{2n}/10\), and then \(\Omega\) large in \eqref{eq:exitlower} and
\eqref{eq:exitupper}; lattice points converging to \(X_i(\delta)\) satisfy
\(W_j(X_0^\Omega)\le\varepsilon\) and lie in \(G_0\) eventually. This proves
\eqref{eq:iterated}. No constant above is controlled uniformly as \(\delta\to0\), which is why
the limit is iterated.\qed

\section{Certificates, replay, and machine-checked algebra}\label{app:certificates}

\subsection{Constants of the operating contract}
All quantities are rational. The chart has \(d=31\) coordinates. Let \(P_i,R_i\) be the saved
matrices of Proposition~\ref{prop:foursinks}, let \(\mathsf c=50000000/10001\) be half of the
largest rate constant, and divide all rate constants by \(\mathsf c\). Then
\begin{itemize}
\item \(\overline p_i=\|P_i\|_\infty\), and
\(\underline p_i=\tfrac9{10}\bigl(\|R_i\|_\infty\|R_i\|_1\bigr)^{-1}\), because
\(R_iP_iR_i^T\ge\tfrac9{10}I\) implies \(P_i\ge\tfrac9{10}(R_i^TR_i)^{-1}\);
\item \(\mu=\tfrac9{10}/\mathsf c\), from the margin of \(-(J^TP_i+P_iJ)\) at the unscaled
rates;
\item \(K=\sum12\,k\,\|\nu\|_1\,g\), summed over the \(24\) binding reactions with constant
\(k\), where \(g=31\) if the substrate reactant is \(S_\varnothing\) and \(g=1\) otherwise.
This bounds the bilinear form of the quadratic part of \(H\): the deviation of \(E\) or \(F\)
is a sum of \(12\) chart coordinates and that of \(S_\varnothing\) a sum of \(31\);
\item \(\varrho_i=\min\{\min_jx^*_j/(4d),\ \mathrm{gap}/16,\ \mu/(4\overline p_iK)\}\), which
keeps the ball in the positive class and the readout, a sum of four chart coordinates, within
\(\mathrm{gap}/8\) of its value at the sink;
\item \(C_i\) by \eqref{eq:localC}; \(v_i\), \(\lambda_i\), \(\trec\), \(T\) and \(\Omega\) as in
Theorem~\ref{thm:finite} and Corollary~\ref{cor:contract}.
\end{itemize}
The integer preparations are verified at \(220\) digits after refining the isolating intervals
by exact bisection. The bounds \(b_G,d_G\) of Proposition~\ref{prop:example-robust} are sums
over reactions of \(k\|\nu\|_1\) times products of the bounds \eqref{eq:localC}, and of the
corresponding derivatives.

\subsection{Files and replay}
The data files are: \path{postproof_candidate_interface.json} (loadings, rate constants,
integer eliminant); \path{postproof_candidate_sinks.json} (isolating intervals, \(P_i\),
\(R_i\)); \path{postproof_candidate_operating.json} and
\path{postproof_candidate_local_certificate.json} (contract constants, \(\Omega\), robustness
thresholds); \path{postproof_candidate_preparations.json} (integer count vectors);
\path{tree_weights.json}; and the corresponding files for the original source. The program
\path{check_paper.py} recomputes, without reading any intermediate result, the tree weights
from \eqref{eq:example-rates} by Laplacian cofactors, the identity \eqref{eq:target}, the
eliminant from \eqref{eq:elim}, the root census, the four interval Lyapunov certificates, all
constants of the previous subsection, the budget \eqref{eq:budget}, the error bounds, the
inventories and preparation inequalities of the four count vectors, the entries of
Table~\ref{tab:resources}, Proposition~\ref{prop:obstruction} and the thresholds of
Proposition~\ref{prop:example-robust}. It also checks symbolically the elimination identity
behind \eqref{eq:elim}, the square weights \eqref{eq:square} and the expansion
\eqref{eq:drift-exact}. It performs no search, optimization, eigenvalue computation or
simulation; every decision is made in exact rational arithmetic or in integer interval
arithmetic with outward rounding. It needs Python with SymPy and runs in about one minute
(\(745\) checks); the option \texttt{-{}-full} (\(776\) checks) also replays the four sink
certificates of the original source. The script \path{diagnostics.py} produces the floating-point diagnostics
quoted in Section~\ref{sec:example}, which are marked as such, and \path{make_figures.py}
the figures. How the certificates were found (high-precision eigenvectors for \(P_i\), a linear
program for the loadings) is irrelevant for their validity.

\subsection{Machine-checked algebra}
Three small Lean~4 files \cite{Lean4,Mathlib} check algebraic identities used in the text:
\texttt{TreeTotals} (the relations \(M=L+(r-u)D\) and \(su L=r-u\), and the positivity of the
derivative used at \(u=r\) in Proposition~\ref{prop:count}); \texttt{UniformLoading} (the
kernel identity of Lemma~\ref{lem:kernel} in scalar form, the cancellation of the zero-order
term and the form of the numerator in \eqref{eq:drift-exact}--\eqref{eq:numerator}, and the
determinant \(-31410\) in Lemma~\ref{lem:rank}); and \texttt{SymmetricAggregation} (the
multiplicity identity of Propositions~\ref{prop:lift} and~\ref{prop:lump} and the scalar form
of \eqref{eq:LV}). They compile without \texttt{sorry} and with warnings treated as errors.
They are routine identities. Sard's theorem, degree theory, the stability constructions,
invariant foliations, the large-deviation estimates and the operating contract are \emph{not}
formalized, and the interval certificates are checked by the Python program above, not by the
Lean kernel.
